\scanpage{13}
\uline{Questions sur courbes elliptiques} sur $\overline{\unsure{Q|\mathbb Q}}$.

Soit $C$ une telle courbe, $G$ son gpe de Galois \unsure{motivique|motiviques} [$G=\mathrm{Gl}(2,Q)$ ou $G=\prod_{K/Q}G_m$, $K$ le corps de multiplication complexe]. Soit $g\in G(\mathbb C)/G(\overline Q)$ définissant l'application du réseau transcendant de Betti dans celui de De Rham.

a) $Q(g)$ de degré de transcendance égal à $\dim\unsure{G|C}$ (4, ou 2) ?

b) Soit $T$ le lieu de $g$ sur $G_{\overline Q}$, \deleted{\unsure{que le}} qu'on peut supposer (par translation à droite) contenir $e$, \deleted{et} soit $H$ le groupe engendré \deleted{\unsure{et}} par $T$. \deleted{et ses conjugués} A-t-on $H=G$

\deleted{\unsure{$g$}\illegible{1}}
c) Soit \deleted{$g'\in H_{\overline Q}\subset G'_{\overline Q}$ le \unsure{sous-groupe} qui conserve} \inserted{$h\in\mathbb C/\overline Q$ [quotient additif)} définissant le splitting transcendant \unsure{(via th.|(i.e. th.} de Hodge) de De Rham. Peut-on avoir $h\in\overline Q$, ou $h$ est-il tjrs transcendant sur $Q$
\margin{\uline{NB} Si $C$ est à multiplication complexe, on a tjrs $h\in\overline Q$ !}
[\uline{cas où} $C$ \uline{n'a pas de multiplication complexe}]
% archive sheet number: 12
